% ===================================================================
%  TABELL Nr 4 — Mogen ålder och ålderdom.
%  Traduction 2026 d'apres texto/io, controlee sur texto/fr et
%  texto/es. Consigne : voir texto/sv/10-tabell-01.tex.
% ===================================================================

%%K t04-apar-1 apar
\VUsaut{4.0mm}
\VUtitre{15.0pt}{60}{TABELL Nr 4}
\VUsaut{1.6mm}
\VUtitre{11.4pt}{0}{Mogen ålder och ålderdom.}
\VUsaut{3.2mm}

%%K t04-tit-1 sub
\VUcentre{11.4pt}{0}{\textit{Första scenen.}}
\VUsaut{1.6mm}
\VUtitre{11.4pt}{0}{Bröllopsmiddagen.}
\VUsaut{2.6mm}

%%K t04-tit-2 sub
\VUpk{10.2pt}{40}{Hösten.}
\VUsaut{3.0mm}

%%K t04-c1-01-1 p
1. --- De unga har vuxit upp. En av dem, Johan, gifter sig. Vi bevistar
\VUgras{bröllopsmiddagen} som samlar makarnas familjer kring samma bord.

%%K t04-c1-02-1 p
2. --- Vi är på \VUgras{hösten}, i \VUgras{september månad}. Den kalla vinden i
\VUgras{oktober} och \VUgras{november} har ännu inte berövat landsbygden dess
gröna lövverk. Kvällsluften är ljum och doftande; därför äger middagen rum under
\VUgras{verandan} \textsuperscript{(8)} vid trädgården.

%%K t04-c1-03-1 p
3. --- Långt bort, på \VUgras{kullen} \textsuperscript{(3)}, ligger Johans föräldrars hem, ett
elegant \VUgras{slott} \textsuperscript{(1)} dolt i grönskan; vackra vingårdar, som ger ett
utmärkt vin, täcker kullens sluttning. Vinodlaren odlar och sköter dem.

%%K t04-c1-04-1 p
4. --- Över vingårdarna flyger en flock \VUgras{rapphöns} \textsuperscript{(2)} som skrämts av
\VUgras{jaktvaktarens} \textsuperscript{(5)} annalkande. Denne genomsöker, med
\VUgras{geväret} under armen och en \VUgras{jaktväska över axeln}, \VUgras{snåret}
\textsuperscript{(4)} med sin \VUgras{hund} \textsuperscript{(6)}. Han övervakar godset och
hindrar tjuvskyttarna från att utrota villebrådet.

%%K t04-c1-05-1 p
5. --- Utanför \VUgras{verandan} klättrar en \VUgras{spaljévinranka} från vilken täta
\VUgras{druvklasar} \textsuperscript{(7)} hänger.

%%K t04-c1-06-1 p
6. --- De vackra höstblommorna, som pryder \VUgras{bordet} \textsuperscript{(16)}, är
\VUgras{krysantemer} \textsuperscript{(9)}; brudens
\VUgras{far} \textsuperscript{(10)} har satt en av dem i knapphålet på sin frack.

%%K t04-c1-07-1 p
7. --- Måltiden lider mot sitt slut. \VUgras{Betjänten} \textsuperscript{(11)} börjar bära in
efterrättsfaten och ska snart ta bort de olika delarna av kuvertet, \VUgras{skedar}
\textsuperscript{(14)}, \VUgras{gafflar} \textsuperscript{(12)}, \VUgras{knivar}
\textsuperscript{(13)}, \VUgras{stora fat} \textsuperscript{(15)} som man inte längre behöver.

%%K t04-tit-3 sub
\VUpk{10.2pt}{40}{Släkten.}
\VUsaut{3.0mm}

%%K t04-c1-08-1 p
8. --- Alla ser glada ut, och leendet med vilket brudgummens \VUgras{mor}
\textsuperscript{(17)} betraktar sina barn bevisar hennes lycka.

%%K t04-c1-09-1 p
9. --- Det giftermålet förenar två rika familjer i trakten. Johan,
\VUgras{maken} \textsuperscript{(18)}, är \VUgras{son} till en stor godsägare
och han blir \VUgras{svärson} till en skicklig industriman.

%%K t04-c1-10-1 p
10. --- \VUgras{Makarna} är unga och ändå var de sedan länge
\VUgras{förlovade}. \VUgras{Den unga kvinnan} \textsuperscript{(19)}, Marta, är
förtjusande i sin \VUgras{vita klänning} prydd med apelsinblommor.

%%K t04-c1-11-1 p
11. --- Hennes \VUgras{svärfar} \textsuperscript{(20)} är mycket lycklig att ha en så älskvärd
\VUgras{svärdotter}, och Johans \VUgras{svärmor} \textsuperscript{(21)} ler när hon ser sin
\VUgras{dotter} så vacker.

%%K t04-c1-12-1 p
12. --- Vid bordets ände sitter de övriga \VUgras{gästerna} \textsuperscript{(22)},
\VUgras{släktingar, vänner, kusiner, kusiner}. En
\VUgras{hovmästare} \textsuperscript{(23)}, med en servett under armen, betjänar
dem flitigt.

%%K t04-tit-4 sub
\VUpk{10.2pt}{40}{Vännerna.}
\VUsaut{3.0mm}

%%K t04-c1-13-1 p
13. --- På bordets högra sida, här är en \VUgras{fröken} \textsuperscript{(24)},
\VUgras{väninna} till den unga bruden, sedan hennes
\VUgras{bror}, som är hennes \VUgras{marskalk} \textsuperscript{(25)}. Han
pratar med sin \VUgras{tärna} \textsuperscript{(26)} syster till brudgummen och som genom det
giftermålet blir Martas \VUgras{svägerska}. Hon är en förtjusande ung flicka. Hon har vackra
\VUgras{smycken, ett pärlhalsband}, och ett
\VUgras{armband} av guld.

%%K t04-c1-14-1 p
14. --- Bredvid dem är herrn som står upp och lyfter sitt glas en \VUgras{vän}
\textsuperscript{(27)} till familjen. Han är en av \VUgras{brudgummens vittnen}. Han
\VUgras{utbringar en skål}, dricker för de unga makarnas
\VUgras{hälsa} och önskar dem en lång lycka. Han är klädd i högtidsdräkt, i frack
med \VUgras{lackskor} \textsuperscript{(28)}. Han är kavaljer åt
\VUgras{mormodern} \textsuperscript{(29)} som är \VUgras{änka}.

%%K t04-c1-15-1 p
15. --- Den unge mannen i \VUgras{frack} \textsuperscript{(31)}, med en tunn svart mustasch,
är Johans \VUgras{svåger} \textsuperscript{(30)}. Han är en framstående ingenjör. Han sitter
bredvid en mycket elegant \VUgras{ung dam} \textsuperscript{(33)}, Martas
\VUgras{äldre syster} och mor till den nätta lilla flickan som kommer, med armen om sin kusin,
för att bjuda en blomma och
\VUgras{önska} honom \VUgras{god afton}. Den damen har mycket vackra
\VUgras{örhängen} och en elegant \VUgras{huvudbonad}.

%%K t04-c1-16-1 p
16. --- Den lille kusinen, så egendomlig i sin \VUgras{blus} \textsuperscript{(36)} och sina
pösiga \VUgras{byxor} \textsuperscript{(37)} är mycket ömkansvärd. Han är en
\VUgras{föräldralös} \textsuperscript{(35)}. Mycket ung förlorade han sin far och sin mor.
Hans farbror är hans \VUgras{förmyndare} \textsuperscript{(38)}, och har förblivit ogift för
att uppfostra det stackars barnet. Han är en godhjärtad man. Han är en aning skallig och
mycket närsynt; han använder en \VUgras{pincené}.

%%K t04-tit-5 sub
\VUsaut{2.4mm}
\VUpk{10.2pt}{40}{Efterrätten.}
\VUsaut{3.0mm}

%%K t04-c1-17-1 p
17. --- Bakom gästerna, på ett bord täckt med en \VUgras{bordduk}, är
\VUgras{efterrätten} \textsuperscript{(39)} framdukad. I en stor skål, här är
frukter, \VUgras{persikor} \textsuperscript{(40)}, \VUgras{päron} \textsuperscript{(41)},
\VUgras{äpplen} \textsuperscript{(42)},
\VUgras{aprikoser} \textsuperscript{(43)} och \VUgras{druvor}
\textsuperscript{(44)}. Nära intill, \VUgras{ananaser} \textsuperscript{(45)}, en vacker
\VUgras{tårta} \textsuperscript{(46)},
\VUgras{likörflaskor} \textsuperscript{(48)} och \VUgras{champagne}
\textsuperscript{(49)} och också staplar av rena \VUgras{tallrikar} \textsuperscript{(47)}.

%%K t04-c1-18-1 p
18. --- \VUgras{Källarmästaren} \textsuperscript{(50)} drar upp en \VUgras{flaska}
\textsuperscript{(52)} gammalt vin med en \VUgras{korkskruv} \textsuperscript{(51)}.
\VUgras{Korken} \textsuperscript{(53)} verkar mycket svår att dra ut. Den källarmästaren har
en \VUgras{servett} \textsuperscript{(54)} under armen.

%%K t04-c1-19-1 p
19. --- Efter middagen ska herrarna gå till rökrummet, och damerna ska gå och göra sig i
ordning till balen.

%%K t04-tit-6 sub
\VUsaut{5.0mm}
\VUcentre{11.4pt}{0}{\textit{Andra scenen.}}
\VUsaut{1.6mm}
\VUpk{11.4pt}{40}{Farfaderns födelsedag.}
\VUsaut{2.6mm}

%%K t04-tit-7 sub
\VUpk{10.2pt}{40}{Besöket.}
\VUsaut{3.0mm}

%%K t04-c2-01-1 p
1. --- Tio år har gått, Johans föräldrar har blivit gamla och älskar mycket sina tre barnbarn,
två \VUgras{sonsöner} \textsuperscript{(2)} och en
\VUgras{sondotter} \textsuperscript{(3)}. Eftersom det i dag är
\VUgras{farfaderns födelsedag}, kommer barnen för att önska dem en god namnsdag.

%%K t04-c2-02-1 p
2. --- Lille Peter är ännu mycket ung, han är bara tre år. Han ser
\VUgras{katten} \textsuperscript{(1)} närma sig. Han vill smeka dess vackra
sidenlena \VUgras{päls} och dess långa \VUgras{svans}, han tänker inte på, att det under de
graciösa \VUgras{tassarna} döljer sig elaka spetsiga klor som kanske ska riva honom.

%%K t04-c2-03-1 p
3. --- Hans syster Gabriella, som ger honom handen, är mycket allvarligare och Marcel med sin
stora \VUgras{kappa} \textsuperscript{(4)} och sitt \VUgras{läderbälte} \textsuperscript{(5)}
ser något manligare ut, trots sina korta byxor som låter se hans \VUgras{strumpor}
\textsuperscript{(6)}.

%%K t04-c2-04-1 p
4. --- Deras far har tagit på sig sin tjocka \VUgras{vinterrock} \textsuperscript{(7)} med
\VUgras{pälskrage} \textsuperscript{(8)} och i sin
\VUgras{ficka} \textsuperscript{(9)} har han en halsduk. Hans hustru har sin
filthatt prydd med \VUgras{fjädrar} \textsuperscript{(10)}, sin \VUgras{jacka}
\textsuperscript{(11)}, sin \VUgras{pälsboa} \textsuperscript{(12)} och sin
\VUgras{muff} \textsuperscript{(13)}.

%%K t04-c2-05-1 p
5. --- Det är mycket kallt, för vintern råder. \VUgras{Snön} \textsuperscript{(19)} har fallit
hela dagen och husens tak är alldeles vita. Samtidigt med \VUgras{natten} blir kylan ännu
skarpare. På himlen lyser
\VUgras{månen} \textsuperscript{(17)} och \VUgras{stjärnorna}
\textsuperscript{(18)} och genomtränger \VUgras{mörkret}.

%%K t04-tit-8 sub
\VUsaut{4.2mm}
\VUpk{10.2pt}{40}{Sjukdomen.}
\VUsaut{3.0mm}

%%K t04-c2-06-1 p
6. --- Den stackars \VUgras{blinde} \textsuperscript{(20)} som går över torget framför
\VUgras{apoteket} \textsuperscript{(16)} ledd av sin \VUgras{pudel} \textsuperscript{(21)}, är
mycket olycklig. Justina, \VUgras{tjänsteflickan} \textsuperscript{(14)}, bär från köket en
\VUgras{skål} \textsuperscript{(15)} mycket varmt \VUgras{örtte} som är rikligt sockrat.

%%K t04-c2-07-1 p
7. --- \VUgras{Mormodern} \textsuperscript{(23)} som vill söka sin
\VUgras{lornjett} \textsuperscript{(26)} på bordet, för att läsa \VUgras{receptet}
\textsuperscript{(27)} som \VUgras{läkaren} \textsuperscript{(25)} just har skrivit, reser sig
ur sin länstol och stöder sig på sin \VUgras{käpp} \textsuperscript{(24)}.

%%K t04-c2-08-1 p
8. --- Ofta lider hon av \VUgras{opasslighet, svindel, förkylningar, tandvärk}, hennes ben är
svaga och hennes ögon är inte längre särskilt goda. Trots det gycklar hon gärna med sin gamle
\VUgras{läkare} som alltid vill skriva ut en
\VUgras{mixtur} \textsuperscript{(28)} åt henne. I dag är hon mycket glad
eftersom hon har sin unga familj omkring sig.

%%K t04-c2-09-1 p
9. --- Det är behagligt i det väl stängda rummet. I den marmorklädda
\VUgras{spisen} \textsuperscript{(29)} flammar en \VUgras{präktig eld}
\textsuperscript{(30)}, och man känner sig lycklig i det milda \VUgras{ljuset} från
\VUgras{lampan} \textsuperscript{(31)} som man just har tänt.

%%K t04-tit-9 sub
\VUsaut{4.2mm}
\VUpk{10.2pt}{40}{Farfadern.}
\VUsaut{3.0mm}

%%K t04-c2-10-1 p
10. --- Till och med \VUgras{farfadern} \textsuperscript{(33)} har glömt, att han var sjuk,
att hans \VUgras{reumatism} fjättrar honom vid hans
\VUgras{länstol} \textsuperscript{(34)} och att han i går ännu hade \VUgras{feber
och svimningar}. Han minns inte längre, att han förra vintern led av
\VUgras{lunginflammation} och att en hes och plågsam \VUgras{hosta} gjorde honom
mycket lidande.

%%K t04-c2-10-2 p
Och ändå blev han mycket svårt botad. Nu mår han något bättre. Men hans ben som han stöder på
en \VUgras{kudde} \textsuperscript{(38)} lagd på en liten
\VUgras{pall} \textsuperscript{(39)} gör också ont på honom.

%%K t04-c2-11-1 p
11. --- När han är omsorgsfullt insvept i sin varma \VUgras{nattrock} \textsuperscript{(35)}
med tjocka \VUgras{knappar} \textsuperscript{(36)} av pärlemor och när han har bredvid sig,
för att läsa tidningen för honom, den lilla \VUgras{föräldralösa flickan}
\textsuperscript{(40)} klädd i en vit
\VUgras{hätta}, i en blå \VUgras{klänning} \textsuperscript{(42)} och en
\VUgras{pelerin} \textsuperscript{(41)}, klagar han inte.

%%K t04-c2-12-1 p
12. --- Varje år, när \VUgras{kalendern} \textsuperscript{(43)} på nytt visar
\VUgras{datumet} den första \VUgras{december}, väntar han otåligt på sina
\VUgras{barnbarn}, för han vet, att de inte ska glömma det viktiga
\VUgras{datumet}.

%%K t04-c2-13-1 p
13. --- Han minns med rörelse den mycket avlägsna tid då han personligen kom och önskade den
ärorike kejserlige \VUgras{marskalken} en god namnsdag, vars
\VUgras{porträtt} \textsuperscript{(44)} hänger på väggen, och som är hans barns
\VUgras{farfars far}.
\VUsaut{6.0mm}
\VUfilet{22mm}
